% GENERATED by nonlinear_modelling/chapter_wide_race.R -- do not hand-edit.
\begin{table}[htbp]\centering
\caption{The shape of each series in Chapter Two, chosen by one rule.}
\label{tab:ch2shapes}
\begin{threeparttable}
\begin{tabular}{p{5.2cm} c l l c}
\toprule
Series & $n$ & Preferred shape & Detail & Flat / linear rejected by \\
\midrule
Fig 2.2  Great office departures & 46 & \textbf{flat} & --- &  0.00 /  1.18 \\
Fig 2.3  Cabinet departures & 46 & \textbf{linear trend} & slope +0.2708/yr &  3.94 /  0.00 \\
Fig 2.6  Electoral volatility & 48 & \textbf{curved (spline)} & edf 8.22 & 70.21 / 53.46 \\
Fig 2.8  Agenda diversity (ENI) & 47 & \textbf{step} & break 2014 & 90.77 /  2.43 \\
Fig 2.8  Agenda churn (KL) & 47 & \textbf{linear trend} & slope -0.0027/yr &  1.96 /  0.00 \\
Fig 2.9  Unexpected movements & 46 & \textbf{step} & break 1985 & 54.33 / 58.16 \\
Fig 2.10 Parliamentary dislocation & 48 & \textbf{step} & break 2016 & 23.31 / 24.86 \\
Fig 2.10 Backbench rebellions & 46 & \textbf{curved (spline)} & edf 8.00 & 75.13 / 45.91 \\
Fig 2.11 Policy agenda churn & 43 & \textbf{flat} & --- &  0.00 /  0.49 \\
Fig 2.12 Machinery of government & 47 & \textbf{linear trend} & slope +0.0289/yr &  0.31 /  0.00 \\
Fig 2.13 NAO delivery failure & 41 & \textbf{step} & break 2015 & 60.71 / 14.63 \\
Fig 2.15 U-turns & 47 & \textbf{flat} & --- &  0.00 /  3.80 \\
Fig 2.16 Turbulence index (composite) & 47 & \textbf{bend (adopted; BIC winner = spline)} & knot 2012 (slope -0.033 $\to$ +0.325) & 35.46 / 32.98 \\
\bottomrule
\end{tabular}
\begin{tablenotes}\footnotesize
\item Six candidate shapes (flat, linear trend, quadratic, step, bend, spline) fitted to each series and compared on BIC, with any break date or knot searched under a 15\% trim and charged as an estimated parameter. The final column gives the BIC penalty carried by the flat and linear-trend models against the preferred one; a penalty above 10 is conventionally decisive, and below 2 means no preference can be claimed.
\item \textbf{Four series prefer a step}, and each lands on the year already reported in the text: agenda diversity 2014, unexpected movements 1985, parliamentary dislocation 2016, NAO delivery failure 2015. The dates are unchanged; what is new is that the step model is now \emph{preferred} rather than merely fitted.
\item \textbf{Backbench rebellions reject a step decisively} ($\Delta$BIC 31.75) and prefer a smooth curve, as does electoral volatility. \textbf{Policy agenda churn and U-turns prefer a flat series} --- no trend and no turning point. Machinery of government prefers a linear trend, but by only 1.30 over flat.
\item Machinery of government, NAO and U-turns are sparse annual counts, for which a Gaussian criterion is indicative only; the Poisson sup-LR run separately on machinery of government agrees ($p=0.10$, no identified break).
\item Source: \texttt{nonlinear\_modelling/chapter\_wide\_race.R}, \texttt{chapter\_wide\_shapes.csv}. Seed 42.
\end{tablenotes}
\end{threeparttable}
\end{table}
